\maketitle
\normalfont
\vspace{-1.2em}
\enlargethispage{2\baselineskip}

\begin{center}
\textbf{摘\quad 要}
\end{center}

\small
霍尔木兹海峡是全球原油运输的关键咽喉。若发生封锁，供应中断、地缘风险重估、战略储备释放、商业库存缓冲、绕道运输恢复和市场预期修复会共同作用于国际原油价格。本文以附件布伦特原油期货主力合约价格数据为基础，建立“数据清洗--传统供需基准--短期动态递推--长期情景预测--敏感性分析”的模型链条。

针对任务一，本文清洗 2017 年 9 月至 2026 年 5 月的原始价格序列，截取 2026 年 3 月 2 日至 2026 年 5 月 5 日为冲突窗口。附件数据显示，窗口内最高收盘价为 114.06 美元/桶，最高盘中价为 119.50 美元/桶。按题面短期低弹性和供应中断范围建立传统供需基准模型，可得 278--337 美元/桶的理论价格区间，显著高于附件真实价格，说明必须引入政策缓冲和市场预期修复。

为解释日度价格路径，本文构建综合机制递推模型：物理供需层纳入供应中断、SPR 释放、商业库存、绕道运输、需求收缩和恐慌衰减，市场价格形成层纳入地缘风险溢价、冲击不确定性、持续封锁制度风险、缓冲确认折价和预期修复折价。综合最优模型在 46 个交易日内取得 RMSE 为 3.44 美元/桶、MAE 为 2.85 美元/桶、MAPE 为 2.86\% 的拟合效果，并优于上一日价格朴素基准和滚动 ARIMA 基准。滞后平移、拐点、机制消融、过拟合压力测试和硬编码审计表明，模型不是简单滞后复制，也没有把 120 美元平台写成代码上限。

针对任务二，本文构建 60--180 天乐观、中性、悲观三种情景路径。新版长期自适应反馈模型显示，中性情景第 180 天价格约为 92.43 美元/桶，悲观情景第 180 天约为 106.83 美元/桶，且仍保留二次跳涨风险。敏感性分析表明，封锁风险衰减速度、地缘风险权重、不确定性与制度风险强度是综合敏感度最高的三个因素；这说明长期结果最依赖市场吸收封锁制度风险的速度，而非某个单一政策释放量。2000 次蒙特卡洛情景树显示，外推期最高价突破 120 美元/桶的条件概率约为 4.6\%。为避免长期中性线被误读为逐日精确预测，本文又引入缓和、维持、升级三状态转移扰动，得到第 180 天 5\%--95\% 条件区间 85.43--93.99 美元/桶，外推期最高价 P95 为 124.61 美元/桶。GPR 滞后审计和 2017--2026 年历史同长度窗口检验进一步说明：2026 冲突窗口具有罕见地缘风险和高波动特征，但外部风险指标只作为审计证据，不作为短期同月拟合因子。

\textbf{关键词：} 霍尔木兹海峡；原油价格；动态递推模型；战略石油储备；情景预测；敏感性分析
\normalfont\normalsize
